% Unified current-result section for the rank-eight manuscript.
% Place after the complete R1768 center-domain reduction.  It uses the
% existing labels thm:research-four-general-spherical and
% prop:creative-r1768-sector.  No additional bibliography key is required.

\Needspace{14\baselineskip}
\subsection{Conjugation parity and the exclusion of R1768}
\label{subsec:r1768-conjugation-exclusion}

The following obstruction uses the full multiplicity vector of a tensor
product.  It needs neither a braiding nor a unitary structure.

\begin{lemma}[Pivotal conjugation parity]
\label{lem:pivotal-conjugation-parity}
Let $\mathcal C$ be a pivotal fusion category over $\mathbb C$, with
pivotal structure $j$.  Suppose an object $V$ has a nondegenerate
selfduality $f:V\to V^*$ satisfying
\[
 f^*\circ j_V=\varepsilon f,\qquad \varepsilon\in\{1,-1\}.
\]
Then $X\otimes V\otimes X^*$ has a nondegenerate selfduality of sign
$\varepsilon$ for every object $X$.  In particular, for selfdual simple
objects $V,S$,
\[
 [X\otimes V\otimes X^*:S]\equiv1\pmod2
 \quad\Longrightarrow\quad
 \nu_2(S)=\nu_2(V).
\]
\end{lemma}
\begin{proof}
Suppress the coherent associativity and tensor-dual comparison maps.
For $W=X\otimes V\otimes X^*$, the isomorphism
\[
 h=j_X\otimes f\otimes\operatorname{id}_{X^*}:W\longrightarrow W^*
\]
satisfies
\[
 h^*\circ j_W
 =j_X\otimes(f^*\circ j_V)
       \otimes(j_X^*\circ j_{X^*})
 =\varepsilon h,
\]
because $j_{X^*}=(j_X^*)^{-1}$.
On the $S$-isotypic summand, write this pairing as a fixed simple
selfduality of sign $\delta=\nu_2(S)$ times a nondegenerate bilinear form
on the multiplicity space.  The latter form has sign
$\varepsilon/\delta$.  When $\varepsilon\ne\delta$, it is alternating,
so its dimension is even.
\end{proof}

\begin{theorem}[General-complex exclusion]
\label{thm:r1768-general-exclusion}
The fusion ring $\sha{R8_3_1768}$ has no categorification by a fusion
category over $\mathbb C$.
\end{theorem}
\begin{proof}
By Theorem~\ref{thm:research-four-general-spherical}, existence of any
complex categorification would imply existence of a pseudo-unitary
categorification with the same based ring.  Work in its canonical
positive pivotal structure.

In the canonical basis $(X_0=1,X_1=g,X_2,\ldots,X_7)$ all simples are
selfdual.  Join two labels $i,k$ whenever some $X$ has odd multiplicity
$[X X_i X^*:X_k]$.  Direct multiplication gives exactly the two connected
components
\[
 \{0,2,3,5\},\qquad \{1,4,6,7\}.
\]
Lemma~\ref{lem:pivotal-conjugation-parity} therefore leaves precisely
the two second-indicator vectors
\begin{equation}\label{eq:r1768-parity-signs}
 (\nu_2(X_0),\ldots,\nu_2(X_7))
 =(1,\varepsilon,1,1,\varepsilon,1,\varepsilon,\varepsilon),
 \qquad\varepsilon=\pm1.
\end{equation}

Both vectors contradict the complete necessary center domains.  For
$\varepsilon=1$, the pointed $C_2$ Fourier cone has an exact rational
Farkas separator.  For $\varepsilon=-1$, the complete primitive-$37$
sector reduction used in Proposition~\ref{prop:creative-r1768-sector}
has five branches,
\[
 U+4V=16,\qquad V=0,1,2,3,4,
\]
and the resulting sign vector has an exact rational Farkas separator
in every branch.  These cones retain the full all-order trace domain
for the rational-ratio columns.  Thus this argument introduces no
cutoff on twist orders.

For each of the six separators $y$, the certificate checks
$y^{\mathsf T}A_j\ge0$ on every allowed column, but
$y^{\mathsf T}b<0$.  This contradicts the necessary equation
$Ax=b$ with $x\ge0$.  The positive-sign domain has $10711$ columns;
each of the five negative-sign domains has $6100$ columns.
Consequently neither vector in~\eqref{eq:r1768-parity-signs} can occur,
contradicting existence of the positive representative and hence of
any complex categorification.
\end{proof}

The complete indicator
calculation in the proof of
Proposition~\ref{prop:creative-r1768-sector} forces
\[
 (\nu_2(X_0),\ldots,\nu_2(X_7))=(1,-1,1,1,1,1,1,1).
\]
But the source tensor gives
\[
 X_2gX_2^*
 =g+2X_2+2X_3+2X_4+2X_5+X_6+X_7.
\]
Its negative-sign selfduality would force even multiplicities for the
positive-indicator constituents $X_6,X_7$, both of which occur once.
The six-certificate proof above independently replays precisely the
parts of the complete center argument needed for the exclusion,
without assuming its final forty-branch indicator conclusion.

The categorical inputs to the theorem are the general-complex existence
reduction and completeness of the necessary center domains, including
the saturated conductor-$37$ argument.  Exact separation verifies
infeasibility of those domains; it does not replace their categorical
necessity proofs.  The replay requires no optimization or associator
search.  The explicit pivotal pairing argument has no additional
algebra-existence, Morita-equivalence, or numerical-indicator
functoriality premise.


The exact source tensor, complete parity graph and six rational separators
are replayed by \path{evidence/frontier/replay/new/r1768/replay_palindromic_center.py}.
The source inventory and categorical dependency statements are collected in
\path{evidence/frontier/README.md}.
